\chapter{Chi Tiết Quy Trình Triển Khai Từng Thành Phần}
\label{ch:deployment}

Toàn bộ quy trình xây dựng hệ thống tuân thủ triết lý \textit{Infrastructure as Code (IaC)} và \textit{GitOps/Automation}, giúp hạ tầng có khả năng tái tạo hoàn toàn tự động, loại bỏ các sai sót từ thao tác cấu hình thủ công.

\section{Khởi tạo hạ tầng AWS bằng Terraform}
Hạ tầng mạng và máy chủ ảo được mô tả bằng mã nguồn Terraform trong thư mục \texttt{terraform/}. Trạng thái hệ thống (\textit{Terraform State}) được lưu trữ tập trung và khóa phiên (\textit{State Locking}) an toàn trên AWS S3 Bucket và DynamoDB:

\begin{lstlisting}[language=bash, caption={Quy trình thực thi khởi tạo hạ tầng Terraform}]
# Khoi tao cac Provider va ket noi S3 Backend
terraform init

# Kiem tra ke hoach thay doi tai nguyen
terraform plan

# Ap dung trien khai ha tang (28 tai nguyen AWS)
terraform apply -auto-approve
\end{lstlisting}

\subsection{Cấu hình Dual-ENI cho máy chủ Edge Gateway}
Nhằm đảm bảo vai trò định tuyến và cô lập mạng, máy chủ Gateway được trang bị 2 Card mạng ảo (\textit{Elastic Network Interface - ENI}):
\begin{lstlisting}[language=bash, caption={Đoạn mã cấu hình Dual-ENI trong Terraform}]
resource "aws_instance" "pfsense" {
  ami           = data.aws_ami.ubuntu.id
  instance_type = "t3.small"

  network_interface {
    network_interface_id = aws_network_interface.pfsense_wan.id # Card ens5 (Public)
    device_index         = 0
  }
  network_interface {
    network_interface_id = aws_network_interface.pfsense_lan.id # Card ens6 (Private)
    device_index         = 1
  }
}
\end{lstlisting}

\section{Cấu hình Edge Gateway \& Hệ thống Suricata IDS}
Được tự động hóa bằng Ansible playbook \texttt{ansible/site.yml} áp dụng role \texttt{gateway}:

\subsection{1. Cấu hình định tuyến hạt nhân và IPTables NAT}
Kích hoạt chuyển tiếp gói tin hạt nhân Linux và thiết lập bảng chuyển dịch cổng:
\begin{lstlisting}[language=bash, caption={Cấu hình IPTables Port Forwarding và SNAT trên Gateway}]
# Kich hoat IP Forwarding
sudo sysctl -w net.ipv4.ip_forward=1

# Chuyen tiep cong 80 vao Ingress-NGINX Master Node
sudo iptables -t nat -A PREROUTING -i ens5 -p tcp --dport 80 -j DNAT --to-destination 10.0.10.10:80

# Chuyen tiep cong 7443 vao Rancher NodePort tren Master Node
sudo iptables -t nat -A PREROUTING -i ens5 -p tcp --dport 7443 -j DNAT --to-destination 10.0.10.10:30443

# Chuyen tiep cong 8443 vao Wazuh Manager REST API
sudo iptables -t nat -A PREROUTING -i ens5 -p tcp --dport 8443 -j DNAT --to-destination 10.0.10.10:55000

# Cau hinh Source NAT Masquerade cho luong Egress ra ngoai Internet
sudo iptables -t nat -A POSTROUTING -o ens5 -j MASQUERADE
sudo iptables -t nat -A POSTROUTING -d 10.0.10.10/32 -o ens6 -j MASQUERADE
\end{lstlisting}

\subsection{2. Cấu hình tập luật phát hiện tấn công cho Suricata}
Suricata được cấu hình lắng nghe trực tiếp trên card WAN \texttt{ens5} ở chế độ \texttt{af-packet}. Tập luật tùy biến tại \texttt{/etc/suricata/rules/local.rules} bao gồm các chữ ký nhận diện tấn công ứng dụng Web:
\begin{lstlisting}[language=bash, caption={Tập luật phát hiện tấn công mạng trong local.rules của Suricata}]
# Nhan dien SQL Injection (UNION SELECT, OR 1=1)
alert http any any -> any any (msg:"ATTACK: SQL Injection Union Select Detected"; flow:to_server,established; content:"UNION"; nocase; content:"SELECT"; nocase; sid:1000001; rev:1;)
alert http any any -> any any (msg:"ATTACK: SQL Injection OR 1=1 Detected"; flow:to_server,established; content:"or"; nocase; content:"1=1"; sid:1000002; rev:1;)

# Nhan dien Directory Traversal / LFI (/etc/passwd, ../..)
alert http any any -> any any (msg:"ATTACK: Directory Traversal /etc/passwd Detected"; flow:to_server,established; content:"/etc/passwd"; sid:1000003; rev:1;)
alert http any any -> any any (msg:"ATTACK: Directory Traversal Path Manipulation"; flow:to_server,established; content:"../.."; sid:1000004; rev:1;)

# Nhan dien Cross-Site Scripting (XSS)
alert http any any -> any any (msg:"ATTACK: XSS Payload Detected"; flow:to_server,established; content:"<script>"; nocase; sid:1000005; rev:1;)

# Nhan dien do quet tep tin he thong nhay cam
alert http any any -> any any (msg:"ATTACK: Sensitive File Probe (.env/.git)"; flow:to_server,established; content:".env"; sid:1000006; rev:1;)
\end{lstlisting}

\section{Khởi tạo cụm Upstream Kubernetes v1.31}
Cụm được triển khai bằng công cụ chính thống \texttt{kubeadm} với phiên bản \textbf{Kubernetes v1.31.14}:
\begin{enumerate}
    \item \textbf{Khởi tạo Control Plane trên Master node}:
    \begin{lstlisting}[language=bash]
sudo kubeadm init --pod-network-cidr=10.244.0.0/16 --apiserver-advertise-address=10.0.10.10
    \end{lstlisting}
    \item \textbf{Cài đặt mạng Flannel CNI}:
    \begin{lstlisting}[language=bash]
kubectl apply -f https://raw.githubusercontent.com/flannel-io/flannel/master/Documentation/kube-flannel.yml
    \end{lstlisting}
    \item \textbf{Gia nhập các Worker node xuyên Vùng sẵn sàng (Multi-AZ)}:
    Chạy lệnh \texttt{kubeadm join} trên Worker 1 (\texttt{10.0.10.139} tại AZ-a) và Worker 2 (\texttt{10.0.11.21} tại AZ-b) để hoàn tất liên kết cụm 3 node.
\end{enumerate}

\section{Triển khai Rancher Management Server}
Rancher Server được triển khai trên cụm K8s thông qua Helm chart chính thức của SUSE:
\begin{lstlisting}[language=bash, caption={Lệnh triển khai Rancher Server bằng Helm}]
helm repo add rancher-prime https://charts.rancher.com/server-charts/prime
helm repo update

# Cai dat Rancher vao namespace cattle-system
helm install rancher rancher-prime/rancher \
  --namespace cattle-system \
  --create-namespace \
  --set hostname=13.251.225.100.sslip.io \
  --set bootstrapPassword=AdminNCSProject2026! \
  --set replicas=1
\end{lstlisting}

Để tách biệt Rancher khỏi cổng Ingress chuẩn 443 của ứng dụng thương mại, một Service loại \texttt{NodePort} chuyên dụng được tạo ra để gắn cổng nội bộ \textbf{\texttt{30443}} với Pod của Rancher:
\begin{lstlisting}[language=yaml, caption={Khai báo NodePort Service cho Rancher Server}]
apiVersion: v1
kind: Service
metadata:
  name: rancher-nodeport
  namespace: cattle-system
spec:
  type: NodePort
  selector:
    app: rancher
  ports:
    - name: https-nodeport
      port: 443
      targetPort: 443
      nodePort: 30443
      protocol: TCP
\end{lstlisting}

\section{Triển khai Ứng dụng Microservices \& Zero-Trust Network Policy}
Toàn bộ 12 vi dịch vụ của Google Online Boutique được triển khai đồng loạt trong namespace \texttt{default}. Để bảo vệ an toàn cho dữ liệu thanh toán và giỏ hàng, chính sách Zero-Trust được áp dụng triệt để:

\begin{lstlisting}[language=yaml, caption={Chính sách NetworkPolicy Zero-Trust bảo vệ PaymentService}]
apiVersion: networking.k8s.io/v1
kind: NetworkPolicy
metadata:
  name: paymentservice-zero-trust
  namespace: default
spec:
  podSelector:
    matchLabels:
      app: paymentservice
  policyTypes:
    - Ingress
    - Egress
  ingress:
    # Chi chap nhan ket noi den tu CheckoutService qua cong gRPC 50051
    - from:
        - podSelector:
            matchLabels:
              app: checkoutservice
      ports:
        - protocol: TCP
          port: 50051
  egress:
    # Chi cho phep Egress giai ma DNS CoreDNS noi bo
    - to:
        - namespaceSelector: {}
          podSelector:
            matchLabels:
              k8s-app: kube-dns
      ports:
        - protocol: UDP
          port: 53
\end{lstlisting}

\section{Cài đặt Wazuh SIEM All-in-One \& Tích hợp Agent}
\begin{enumerate}
    \item \textbf{Cài đặt Wazuh Manager trên Master node}: Cài đặt gói \texttt{wazuh-manager=4.14.9-1}. Tắt module \texttt{<vulnerability-detection>} trong \texttt{/var/ossec/etc/ossec.conf} để tiết kiệm 8.4 GB dung lượng đĩa, bảo đảm máy chủ Master luôn duy trì dung lượng đĩa trống an toàn (>40\%).
    \item \textbf{Cài đặt và kết nối Wazuh Agent trên Gateway}:
    \begin{itemize}
        \item Đăng ký Agent với Manager: \texttt{/var/ossec/bin/agent-auth -m 10.0.10.10 -A pfsense-gateway}.
        \item Cấu hình Agent đọc trực tiếp file nhật ký cảnh báo của Suricata dạng JSON:
        \begin{lstlisting}[language=xml]
<localfile>
  <log_format>json</log_format>
  <location>/var/log/suricata/eve.json</location>
</localfile>
        \end{lstlisting}
    \end{itemize}
\end{enumerate}

\section{Gia cố bảo mật hạ tầng Cloud \& Kubernetes Hardening}
\begin{itemize}
    \item \textbf{Bảo vệ thông tin xác thực AWS qua IMDSv2}: Toàn bộ máy chủ EC2 đều cấu hình \texttt{http\_tokens = "required"} và \texttt{http\_put\_response\_hop\_limit = 2}, ngăn chặn triệt để hình thức tấn công SSRF đánh cắp IAM Role credentials.
    \item \textbf{Mã hóa đĩa cứng tầng lưu trữ}: 100\% các ổ cứng EBS (Root Volume) đều được mã hóa bằng khóa quản lý AWS KMS (\texttt{encrypted = true}).
    \item \textbf{Tính sẵn sàng cao của Pod (PodDisruptionBudget)}: Áp dụng \texttt{minAvailable: 1} cho các dịch vụ cốt lõi (\texttt{frontend}, \texttt{cartservice}, \texttt{checkoutservice}), ngăn ngừa tình trạng gián đoạn khi bảo trì hoặc nâng cấp Node.
    \item \textbf{Chống gom cụm Pod (PodAntiAffinity)}: Cấu hình quy tắc phân tán mềm (\texttt{preferredDuringSchedulingIgnoredDuringExecution}) cho các Pod \texttt{frontend}, đảm bảo các bản sao Pod được ưu tiên xếp đều trên cả 2 Worker node ở 2 Availability Zones khác nhau.
\end{itemize}
